% =====================================================================
%  DANH MỤC CHUẨN — tên bộ tài liệu, tên chương, tên chuyên đề.
%  Đầu trang và tiêu đề mọi tài liệu đều lấy tên từ đây. Đổi tên thì chỉ sửa ở file này.
%  Dùng trong file .tex:  \loai{1}\chuong{3}   hoặc   \loai{2}\chuyende{1}   hoặc   \loai{3}\de{01}
% =====================================================================

% ---------- Tên 3 bộ tài liệu (bên phải đầu trang) ----------
\dmuc{loai1}{CÁC CHUYÊN ĐỀ CƠ BẢN TOÁN 9}
\dmuc{loai2}{CÁC CHUYÊN ĐỀ TOÁN ÔN THI VÀO LỚP 10}
\dmuc{loai3}{LUYỆN ĐỀ THI THỬ VÀO LỚP 10}

% ---------- Loại 1: các chương SGK Toán 9 Kết nối tri thức ----------
\dmuc{chuong1}{Chương I. Phương trình và hệ hai phương trình bậc nhất hai ẩn}
\dmuc{chuong2}{Chương II. Phương trình và bất phương trình bậc nhất một ẩn}
\dmuc{chuong3}{Chương III. Căn bậc hai và căn bậc ba}
\dmuc{chuong4}{Chương IV. Hệ thức lượng trong tam giác vuông}
\dmuc{chuong5}{Chương V. Đường tròn}
\dmuc{chuong6}{Chương VI. Hàm số $y=ax^2\ (a\ne0)$. Phương trình bậc hai một ẩn}
\dmuc{chuong7}{Chương VII. Tần số và tần số tương đối}
\dmuc{chuong8}{Chương VIII. Xác suất của biến cố trong một số mô hình xác suất đơn giản}
\dmuc{chuong9}{Chương IX. Đường tròn ngoại tiếp và đường tròn nội tiếp}
\dmuc{chuong10}{Chương X. Một số hình khối trong thực tiễn}

% Tên ngắn cho đầu trang (chỉ cần với tên dài; tiêu đề trong tài liệu vẫn dùng tên đầy đủ)
\dmucngan{chuong1}{Chương I. Phương trình và hệ hai phương trình bậc nhất hai ẩn}
\dmucngan{chuong2}{Chương II. Phương trình, bất phương trình bậc nhất một ẩn}
\dmucngan{chuong6}{Chương VI. Hàm số $y=ax^2$. Phương trình bậc hai một ẩn}
\dmucngan{chuong8}{Chương VIII. Xác suất của biến cố}
\dmucngan{chuong9}{Chương IX. Đường tròn ngoại tiếp, đường tròn nội tiếp}

% ---------- Loại 2: chuyên đề ôn thi vào 10 Sở Hà Nội ----------
% ĐÃ CHỐT 29/10/2026, khớp 00_danh-muc/dang-CD.md: CD<n> = chuyende<n>, thứ tự theo vị trí câu trong đề Sở HN.
\dmuc{chuyende1}{Chuyên đề 1. Thống kê và xác suất}
\dmuc{chuyende2}{Chuyên đề 2. Biểu thức chứa căn và các câu hỏi phụ}
\dmuc{chuyende3}{Chuyên đề 3. Giải bài toán bằng cách lập phương trình, hệ phương trình, bất phương trình}
\dmuc{chuyende4}{Chuyên đề 4. Phương trình bậc hai, hệ thức Vi-ét, parabol}
\dmuc{chuyende5}{Chuyên đề 5. Hình học thực tế}
\dmuc{chuyende6}{Chuyên đề 6. Hình học phẳng: đường tròn}
\dmuc{chuyende7}{Chuyên đề 7. Bất đẳng thức, cực trị và bài toán tối ưu}

\dmucngan{chuyende3}{Chuyên đề 3. Giải bài toán bằng cách lập PT, hệ PT, BPT}
\dmucngan{chuyende7}{Chuyên đề 7. Cực trị và bài toán tối ưu}